\documentclass[aspectratio=169,10pt]{beamer}

\usepackage{luatexja}
\usepackage{luatexja-fontspec}
\usepackage{amsmath}
\usepackage{amssymb}
\usepackage{mathtools}
\usepackage{bm}
\usepackage{booktabs}
\usepackage{listings}

%==================================================
% 色と見た目
\definecolor{main}{RGB}{30,70,140}
\definecolor{alertcol}{RGB}{200,80,40}
\definecolor{examplecol}{RGB}{30,120,60}
\definecolor{codebg}{RGB}{245,247,250}

\setbeamercolor{structure}{fg=main}
\setbeamercolor{frametitle}{fg=white,bg=main}
\setbeamercolor{block title}{fg=white,bg=main}
\setbeamercolor{block body}{fg=black,bg=main!6}
\setbeamercolor{block title alerted}{fg=white,bg=alertcol}
\setbeamercolor{block body alerted}{fg=black,bg=alertcol!6}
\setbeamercolor{block title example}{fg=white,bg=examplecol}
\setbeamercolor{block body example}{fg=black,bg=examplecol!6}
\setbeamertemplate{blocks}[rounded][shadow=false]
\setbeamertemplate{itemize items}[circle]
\setbeamertemplate{navigation symbols}{}
\setbeamertemplate{footline}[frame number]

\setbeamerfont{block title}{series=\bfseries}
\setbeamerfont{block title alerted}{series=\bfseries}
\setbeamerfont{block title example}{series=\bfseries}

%==================================================
% コード表示
\lstset{
  language=[LaTeX]TeX,
  basicstyle=\ttfamily\scriptsize,
  backgroundcolor=\color{codebg},
  frame=single,
  rulecolor=\color{black!15},
  columns=fullflexible,
  keepspaces=true,
  showstringspaces=false,
  breaklines=true,
  aboveskip=0.4em,
  belowskip=0.4em
}

%==================================================
% よく使うコマンド
\newcommand{\R}{\mathbb{R}}
\newcommand{\E}{\mathbb{E}}
\newcommand{\KL}{D_{\mathrm{KL}}}
\newcommand{\data}{p_{\mathrm{data}}}
\newcommand{\Normal}{\mathcal{N}}
\DeclareMathOperator*{\argmin}{arg\,min}
\DeclareMathOperator*{\argmax}{arg\,max}
\DeclareMathOperator{\Tr}{Tr}
\DeclareMathOperator{\diag}{diag}

\title{機械学習・拡散モデルのための\\LaTeX / Beamer 早見表}
\subtitle{数式・ブロック・箇条書き・色・レイアウト}
\author{ゼミ資料作成用}
\date{\today}

\begin{document}

\begin{frame}
  \titlepage
\end{frame}

\begin{frame}{この早見表の使い方}
  \begin{block}{目的}
    左側のコードをコピーし、右側の表示を確認しながらスライドを作る。
  \end{block}

  \begin{itemize}
    \item 数式は『The Principles of Diffusion Models』の記法を中心に収録
    \item ソース：\texttt{tex/latex-cheatsheet.tex}
    \item 完成PDF：\texttt{pdf/latex-cheatsheet.pdf}
    \item 新しいスライドは基本的に
      \texttt{\textbackslash begin\{frame\}\{タイトル\}} から作る
  \end{itemize}

  \begin{alertblock}{注意}
    \texttt{\textbackslash begin} と \texttt{\textbackslash end}、
    \texttt{\{} と \texttt{\}} は必ず対応させる。
  \end{alertblock}
\end{frame}

\section{基本}

\begin{frame}[fragile]{最小構成}
\begin{lstlisting}
\documentclass[aspectratio=169]{beamer}
\usepackage{luatexja}
\usepackage{luatexja-fontspec}

\title{タイトル}
\author{著者名}
\date{\today}

\begin{document}
\begin{frame}
  \titlepage
\end {frame}

\begin{frame}{スライドタイトル}
  本文
\end {frame}
\end {document}
\end{lstlisting}

\begin{exampleblock}{ポイント}
  日本語をLuaLaTeXで扱うため、\texttt{luatexja}を読み込む。
  \texttt{aspectratio=169}は16:9のスライドを意味する。
\end{exampleblock}
\end{frame}

\begin{frame}[fragile]{文字と数式の基本}
\begin{columns}[T]
\begin{column}{0.50\textwidth}
\begin{lstlisting}
本文中の数式は $p(x)$ と書く。

独立した数式は
\[
  p_\phi(x) \approx p_{\mathrm{data}}(x)
\]

\textbf{太字}
\emph{強調}
\alert{注意}
\end{lstlisting}
\end{column}
\begin{column}{0.46\textwidth}
本文中の数式は $p(x)$ と書く。

独立した数式は
\[
  p_\phi(x) \approx p_{\mathrm{data}}(x)
\]

\textbf{太字}\\
\emph{強調}\\
\alert{注意}
\end{column}
\end{columns}

\begin{alertblock}{使い分け}
  数式中の添字は \texttt{\_}、上付きは \texttt{\^{}}。
  本文中の特殊文字 \texttt{\_ \% \& \#} には先頭に
  \texttt{\textbackslash}を付ける。
\end{alertblock}
\end{frame}

\section{機械学習の数式}

\begin{frame}[fragile]{スカラー・ベクトル・行列}
\begin{columns}[T]
\begin{column}{0.52\textwidth}
\begin{lstlisting}
$a$                       % scalar
$\mathbf{x}\in\mathbb{R}^D$
$\bm{x}$                  % bold symbol
$A\in\mathbb{R}^{m\times n}$
$A^\top$
$I_D$
$\operatorname{diag}(\mathbf{a})$
$\operatorname{Tr}(A)$
$\|\mathbf{x}\|_2$
\end{lstlisting}
\end{column}
\begin{column}{0.44\textwidth}
\begin{align*}
a & &&\text{スカラー}\\
\mathbf{x}&\in\mathbb{R}^D &&\text{ベクトル}\\
\bm{x}& &&\text{太字の数式記号}\\
A&\in\mathbb{R}^{m\times n} &&\text{行列}\\
A^\top& &&\text{転置}\\
I_D& &&\text{単位行列}\\
\diag(\mathbf{a})& &&\text{対角行列}\\
\Tr(A)& &&\text{トレース}\\
\|\mathbf{x}\|_2& &&\text{$L_2$ノルム}
\end{align*}
\end{column}
\end{columns}
\end{frame}

\begin{frame}[fragile]{確率・統計}
\begin{columns}[T]
\begin{column}{0.52\textwidth}
\begin{lstlisting}
$\mathbf{x}\sim p_{\mathrm{data}}$
$p(\mathbf{x}\mid\mathbf{z})$
$\mathbb{E}_{\mathbf{x}\sim p}[f(\mathbf{x})]$
$\operatorname{Var}(X)$
$\operatorname{Cov}(X,Y)$
$\epsilon\sim\mathcal{N}(0,I)$
$\mathcal{N}(\mathbf{x};\bm{\mu},\Sigma)$
\end{lstlisting}
\end{column}
\begin{column}{0.44\textwidth}
\begin{align*}
\mathbf{x}&\sim \data\\
p(\mathbf{x}\mid\mathbf{z})&\\
\E_{\mathbf{x}\sim p}[f(\mathbf{x})]&\\
\operatorname{Var}(X)&\\
\operatorname{Cov}(X,Y)&\\
\epsilon&\sim\Normal(0,I)\\
\Normal(\mathbf{x};\bm{\mu},\Sigma)&
\end{align*}
\end{column}
\end{columns}

\begin{block}{読み方}
  \texttt{\textbackslash mid} は条件「given」、
  \texttt{\textbackslash sim} は「～に従う」を表す。
\end{block}
\end{frame}

\begin{frame}[fragile]{微分・勾配・積分}
\begin{columns}[T]
\begin{column}{0.52\textwidth}
\begin{lstlisting}
$\frac{dy}{dx}$
$\frac{\partial f}{\partial x_i}$
$\nabla_{\mathbf{x}} f(\mathbf{x})$
$\nabla_{\mathbf{x}}\log p(\mathbf{x})$
$\nabla\cdot\mathbf{v}(\mathbf{x},t)$
$\nabla_{\mathbf{x}}^2 f(\mathbf{x})$
$\int_{\mathbb{R}^D} f(\mathbf{x})\,d\mathbf{x}$
\end{lstlisting}
\end{column}
\begin{column}{0.44\textwidth}
\begin{align*}
\frac{dy}{dx}& &&\text{常微分}\\
\frac{\partial f}{\partial x_i}& &&\text{偏微分}\\
\nabla_{\mathbf{x}} f(\mathbf{x})& &&\text{勾配}\\
\nabla_{\mathbf{x}}\log p(\mathbf{x})& &&\text{スコア}\\
\nabla\cdot\mathbf{v}(\mathbf{x},t)& &&\text{発散}\\
\nabla_{\mathbf{x}}^2 f(\mathbf{x})& &&\text{ラプラシアン}\\
\int_{\R^D} f(\mathbf{x})\,d\mathbf{x}&
\end{align*}
\end{column}
\end{columns}
\end{frame}

\begin{frame}[fragile]{最適化・損失関数}
\begin{lstlisting}
\[
  \phi^* \in \argmin_{\phi}
  \mathcal{D}(p_{\mathrm{data}},p_\phi)
\]

\[
  \widehat{\mathcal{L}}(\phi)
  = -\frac{1}{N}\sum_{i=1}^{N}
    \log p_\phi\!\left(\mathbf{x}^{(i)}\right)
\]
\end{lstlisting}

\begin{columns}[T]
\begin{column}{0.48\textwidth}
\begin{block}{母集団上の目的}
\[
\phi^*\in\argmin_\phi
\mathcal{D}(\data,p_\phi)
\]
\end{block}
\end{column}
\begin{column}{0.48\textwidth}
\begin{exampleblock}{経験損失の例}
\[
\widehat{\mathcal{L}}(\phi)
=-\frac{1}{N}\sum_{i=1}^{N}
\log p_\phi(\mathbf{x}^{(i)})
\]
\end{exampleblock}
\end{column}
\end{columns}
\end{frame}

\begin{frame}[fragile]{KLダイバージェンス}
\begin{lstlisting}
\begin{align}
D_{\mathrm{KL}}(p\|q)
&= \int p(\mathbf{x})
   \log\frac{p(\mathbf{x})}{q(\mathbf{x})}
   \,d\mathbf{x} \\
&= \mathbb{E}_{\mathbf{x}\sim p}
   [\log p(\mathbf{x})-\log q(\mathbf{x})].
\end{align}
\end{lstlisting}

\begin{block}{表示}
\begin{align*}
\KL(p\|q)
&=\int p(\mathbf{x})
\log\frac{p(\mathbf{x})}{q(\mathbf{x})}\,d\mathbf{x}\\
&=\E_{\mathbf{x}\sim p}
[\log p(\mathbf{x})-\log q(\mathbf{x})].
\end{align*}
\end{block}

\begin{alertblock}{重要}
  KLは一般に非対称：
  $\KL(p\|q)\neq\KL(q\|p)$。
  二本線は \texttt{\textbackslash|} ではなく \texttt{\textbackslash\textbackslash|} または
  \texttt{\textbackslash mid\textbackslash mid} より、
  \texttt{\textbackslash Vert} や \texttt{\textbackslash\textbackslash|} の見た目を確認して使う。
\end{alertblock}
\end{frame}

\begin{frame}[fragile]{拡散モデル：ガウス摂動}
\begin{lstlisting}
\[
  p_t(\mathbf{x}_t\mid\mathbf{x}_0)
  = \mathcal{N}
    (\mathbf{x}_t;
     \alpha_t\mathbf{x}_0,\sigma_t^2 I)
\]
\[
  \mathbf{x}_t
  = \alpha_t\mathbf{x}_0+\sigma_t\bm{\epsilon},
  \qquad \bm{\epsilon}\sim\mathcal{N}(0,I).
\]
\end{lstlisting}

\begin{exampleblock}{表示}
\[
p_t(\mathbf{x}_t\mid\mathbf{x}_0)
=\Normal(\mathbf{x}_t;
\alpha_t\mathbf{x}_0,\sigma_t^2I)
\]
\[
\mathbf{x}_t
=\alpha_t\mathbf{x}_0+\sigma_t\bm{\epsilon},
\qquad \bm{\epsilon}\sim\Normal(0,I).
\]
\end{exampleblock}
\end{frame}

\begin{frame}[fragile]{拡散モデル：Score・SDE・ODE}
\begin{lstlisting}
% score function
\[
  \mathbf{s}_\theta(\mathbf{x},t)
  \approx \nabla_{\mathbf{x}}\log p_t(\mathbf{x})
\]

% forward SDE
\[
  d\mathbf{x}_t
  = \mathbf{f}(\mathbf{x}_t,t)\,dt
  + g(t)\,d\mathbf{w}_t
\]

% probability-flow ODE
\[
  \frac{d\mathbf{x}_t}{dt}
  = \mathbf{v}_\theta(\mathbf{x}_t,t)
\]
\end{lstlisting}

\[
\mathbf{s}_\theta(\mathbf{x},t)
\approx\nabla_{\mathbf{x}}\log p_t(\mathbf{x}),
\quad
d\mathbf{x}_t=\mathbf{f}(\mathbf{x}_t,t)\,dt+g(t)\,d\mathbf{w}_t
\]
\[
\frac{d\mathbf{x}_t}{dt}
=\mathbf{v}_\theta(\mathbf{x}_t,t)
\]
\end{frame}

\begin{frame}[fragile]{長い数式を揃える}
\begin{columns}[T]
\begin{column}{0.50\textwidth}
\begin{lstlisting}
\begin{align}
\mathcal{L}(\theta)
&= \mathbb{E}_{t,\mathbf{x}_0,\epsilon}
   \Bigl[
   \lambda(t)
   \bigl\|
   \mathbf{s}_\theta(\mathbf{x}_t,t)
   - \nabla_{\mathbf{x}_t}
     \log p_t(\mathbf{x}_t)
   \bigr\|_2^2
   \Bigr] \\
&\ge 0.
\end{align}
\end{lstlisting}
\end{column}
\begin{column}{0.46\textwidth}
\begin{align*}
\mathcal{L}(\theta)
&=\E_{t,\mathbf{x}_0,\epsilon}
\Bigl[
\lambda(t)
\bigl\|
\mathbf{s}_\theta(\mathbf{x}_t,t)\\
&\qquad-\nabla_{\mathbf{x}_t}
\log p_t(\mathbf{x}_t)
\bigr\|_2^2
\Bigr]\\
&\ge 0.
\end{align*}
\end{column}
\end{columns}

\begin{block}{コツ}
  \texttt{\&}を揃えたい位置に置き、
  \texttt{\textbackslash\textbackslash}で改行する。
  番号不要なら \texttt{align*} を使う。
\end{block}
\end{frame}

\begin{frame}[fragile]{括弧・場合分け・集合}
\begin{columns}[T]
\begin{column}{0.52\textwidth}
\begin{lstlisting}
\[
\left\{\mathbf{x}^{(i)}\right\}_{i=1}^{N}
\]

\[
f(x)=
\begin{cases}
  x^2, & x\ge 0,\\
  0,   & x<0.
\end{cases}
\]

\[
\underbrace{p(\mathbf{x})}_{\text{target}}
\quad
\overbrace{q_\phi(\mathbf{x})}^{\text{model}}
\]
\end{lstlisting}
\end{column}
\begin{column}{0.44\textwidth}
\[
\left\{\mathbf{x}^{(i)}\right\}_{i=1}^{N}
\]
\[
f(x)=
\begin{cases}
x^2,&x\ge0,\\
0,&x<0.
\end{cases}
\]
\[
\underbrace{p(\mathbf{x})}_{\text{target}}
\quad
\overbrace{q_\phi(\mathbf{x})}^{\text{model}}
\]
\end{column}
\end{columns}
\end{frame}

\section{Beamerの機能}

\begin{frame}[fragile]{3種類のブロック}
\begin{columns}[T]
\begin{column}{0.49\textwidth}
\begin{lstlisting}
\begin{block}{定義}
  通常の説明や定義
\end{block}

\begin{exampleblock}{例・結論}
  例、利点、重要な結論
\end{exampleblock}

\begin{alertblock}{注意}
  注意点、制約、誤りやすい点
\end{alertblock}
\end{lstlisting}
\end{column}
\begin{column}{0.47\textwidth}
\begin{block}{定義}
通常の説明や定義
\end{block}
\begin{exampleblock}{例・結論}
例、利点、重要な結論
\end{exampleblock}
\begin{alertblock}{注意}
注意点、制約、誤りやすい点
\end{alertblock}
\end{column}
\end{columns}
\end{frame}

\begin{frame}[fragile]{ブロックの色と形}
\begin{lstlisting}
\definecolor{main}{RGB}{30,70,140}
\definecolor{alertcol}{RGB}{200,80,40}

\setbeamertemplate{blocks}[rounded][shadow=false]
\setbeamercolor{block title}{fg=white,bg=main}
\setbeamercolor{block body}{fg=black,bg=main!6}
\setbeamercolor{block title alerted}
  {fg=white,bg=alertcol}
\setbeamercolor{block body alerted}
  {fg=black,bg=alertcol!6}
\end{lstlisting}

\begin{itemize}
  \item \texttt{fg}：文字色、\texttt{bg}：背景色
  \item \texttt{main!6}：mainを6\%使用した薄い色
  \item \texttt{RGB}：各成分を0～255で指定
  \item \texttt{rounded}：角丸、\texttt{shadow=false}：影なし
\end{itemize}
\end{frame}

\begin{frame}[fragile]{箇条書き・番号付きリスト}
\begin{columns}[T]
\begin{column}{0.50\textwidth}
\begin{lstlisting}
\begin{itemize}
  \item 1段目
  \begin{itemize}
    \item 2段目
  \end{itemize}
  \item \alert{重要項目}
\end{itemize}

\begin{enumerate}
  \item forward process
  \item training
  \item sampling
\end{enumerate}
\end{lstlisting}
\end{column}
\begin{column}{0.46\textwidth}
\begin{itemize}
  \item 1段目
  \begin{itemize}
    \item 2段目
  \end{itemize}
  \item \alert{重要項目}
\end{itemize}

\begin{enumerate}
  \item forward process
  \item training
  \item sampling
\end{enumerate}
\end{column}
\end{columns}

\begin{block}{記号の形}
\texttt{\textbackslash setbeamertemplate\{itemize items\}[circle]}
の \texttt{circle} は \texttt{square} や \texttt{triangle} に変更できる。
\end{block}
\end{frame}

\begin{frame}[fragile]{2列レイアウト}
\begin{lstlisting}
\begin{columns}[T]
  \begin{column}{0.48\textwidth}
    左側の内容
  \end{column}
  \begin{column}{0.48\textwidth}
    右側の内容
  \end{column}
\end{columns}
\end{lstlisting}

\begin{columns}[T]
\begin{column}{0.48\textwidth}
\begin{block}{左：理論}
\[
\mathbf{x}_t
=\alpha_t\mathbf{x}_0+\sigma_t\epsilon
\]
\begin{itemize}
\item forward process
\item noise addition
\end{itemize}
\end{block}
\end{column}
\begin{column}{0.48\textwidth}
\begin{exampleblock}{右：意味}
データに時間依存のノイズを加え、
単純な事前分布へ徐々に変換する。
\end{exampleblock}
\end{column}
\end{columns}

\begin{alertblock}{幅}
  2列なら各 \texttt{0.48\textbackslash textwidth} 程度にすると余白が残る。
\end{alertblock}
\end{frame}

\begin{frame}[fragile]{図の挿入}
\begin{lstlisting}
\usepackage{graphicx} % beamerでは通常読み込み済み

\begin{figure}
  \centering
  \includegraphics[
    width=0.75\textwidth
  ]{images/figure-name.png}
  \caption{図の説明}
\end{figure}
\end{lstlisting}

\begin{itemize}
  \item 画像はプロジェクト内の \texttt{tex/images/} などにまとめる
  \item \texttt{width}で横幅、\texttt{height}で高さを指定
  \item 縦横比を保つには片方だけ指定する
  \item 論文の図を使う場合は出典と利用条件を確認する
\end{itemize}

\begin{block}{出典の例}
  \texttt{\textbackslash footnotesize Source: Author et al., year, Fig.~1.}
\end{block}
\end{frame}

\begin{frame}[fragile]{表の作成}
\begin{columns}[T]
\begin{column}{0.55\textwidth}
\begin{lstlisting}
\begin{tabular}{lll}
\toprule
View & Object & Dynamics \\
\midrule
Variational & denoiser & discrete \\
Score-based & score & SDE \\
Flow-based & velocity & ODE \\
\bottomrule
\end{tabular}
\end{lstlisting}
\end{column}
\begin{column}{0.41\textwidth}
\centering
\begin{tabular}{lll}
\toprule
View & Object & Eq.\\
\midrule
VAE & denoiser & discrete\\
Score & score & SDE\\
Flow & velocity & ODE\\
\bottomrule
\end{tabular}
\end{column}
\end{columns}

\begin{itemize}
  \item \texttt{l / c / r}：左寄せ・中央・右寄せ
  \item \texttt{\&}：列の区切り、\texttt{\textbackslash\textbackslash}：行の終端
  \item \texttt{booktabs}の横線は見やすい表に適する
\end{itemize}
\end{frame}

\begin{frame}[fragile]{段階表示（発表用）}
\begin{columns}[T]
\begin{column}{0.50\textwidth}
\begin{lstlisting}
\begin{itemize}
  \item<1-> 最初に表示
  \item<2-> 次に表示
  \item<3-> 最後に表示
\end{itemize}

\only<2>{
  2枚目だけに表示
}

\uncover<3->{
  3枚目以降に表示
}
\end{lstlisting}
\end{column}
\begin{column}{0.46\textwidth}
\begin{itemize}
  \item<1-> 最初に表示
  \item<2-> 次に表示
  \item<3-> 最後に表示
\end{itemize}

\begin{block}{簡単な方法}
  項目間に \texttt{\textbackslash pause} を置いても段階表示できる。
\end{block}
\end{column}
\end{columns}
\end{frame}

\begin{frame}[fragile]{数式番号と参照}
\begin{lstlisting}
\begin{equation}
  \phi^*
  = \argmin_\phi \mathcal{L}(\phi)
  \label{eq:training}
\end{equation}

式~\eqref{eq:training}より、最適な
パラメータを求める。
\end{lstlisting}

\begin{block}{ルール}
\begin{itemize}
  \item \texttt{\textbackslash label\{eq:名前\}}：番号に名前を付ける
  \item \texttt{\textbackslash eqref\{eq:名前\}}：括弧付きで式番号を参照
  \item 図は \texttt{fig:}、表は \texttt{tab:} を接頭辞にすると整理しやすい
  \item 参照番号が \texttt{??} のときは、通常はもう一度ビルドする
\end{itemize}
\end{block}
\end{frame}

\section{実用上の注意}

\begin{frame}[fragile]{頻出エラーの確認表}
\small
\begin{tabular}{p{0.28\textwidth}p{0.31\textwidth}p{0.33\textwidth}}
\toprule
症状 & 主な原因 & 確認すること\\
\midrule
\texttt{Undefined control sequence}
& コマンド名の誤記、パッケージ不足
& エラー行の \texttt{\textbackslash command} を確認\\
\addlinespace
\texttt{Missing \$ inserted}
& 数式コマンドを本文で使用
& \texttt{\$...\$} 内に入れる\\
\addlinespace
\texttt{Extra \} / runaway argument}
& 波括弧の数が不一致
& \texttt{\{} と \texttt{\}} を対応させる\\
\addlinespace
スライドから内容がはみ出す
& 1枚の情報量が多い
& 分割、\texttt{\textbackslash small}、余白調整\\
\addlinespace
画像が見つからない
& パスまたは拡張子が違う
& texファイルから見た相対パスを確認\\
\bottomrule
\end{tabular}

\begin{alertblock}{基本方針}
  文字を極端に小さくする前に、1枚を2枚へ分割する。
\end{alertblock}
\end{frame}

\begin{frame}{論文解説スライドの推奨構成}
\begin{enumerate}
  \item \textbf{問い}：この節は何を解決するのか
  \item \textbf{前提・記号}：変数と分布を先に定義
  \item \textbf{中心式}：1枚につき主役となる式を1つ
  \item \textbf{式の意味}：各項を文章または色で対応付け
  \item \textbf{直感・図}：forward / reverseなどの流れを図示
  \item \textbf{要点}：結論と次節への接続
\end{enumerate}

\begin{exampleblock}{1枚の基本形}
  上：結論または中心式 \quad
  左下：定義・導出 \quad
  右下：直感・注意点
\end{exampleblock}
\end{frame}

\begin{frame}{論文で特によく使う記号}
\small
\begin{tabular}{lll}
\toprule
記号 & LaTeX & 意味\\
\midrule
$\data$ & \texttt{p\_\{\textbackslash mathrm\{data\}\}} & データ分布\\
$p_\phi$ & \texttt{p\_\textbackslash phi} & モデル分布\\
$\phi^*$ & \texttt{\textbackslash phi\^{}*} & 最適パラメータ\\
$\mathbf{x}\sim p$ & \texttt{\textbackslash mathbf\{x\}\textbackslash sim p} & 分布$p$からの標本\\
$\E[f(\mathbf{x})]$ & \texttt{\textbackslash mathbb\{E\}[f(\textbackslash mathbf\{x\})]} & 期待値\\
$\KL(p\|q)$ & \texttt{D\_\{\textbackslash mathrm\{KL\}\}(p\textbackslash|q)} & KL divergence\\
$\nabla_{\mathbf{x}}\log p_t$ & \texttt{\textbackslash nabla\_\{\textbackslash mathbf\{x\}\}\textbackslash log p\_t} & score\\
$\Normal(0,I)$ & \texttt{\textbackslash mathcal\{N\}(0,I)} & 標準正規分布\\
$\mathbf{x}_t$ & \texttt{\textbackslash mathbf\{x\}\_t} & 時刻$t$の状態\\
\bottomrule
\end{tabular}
\end{frame}

\begin{frame}{参考資料}
\begin{block}{主資料}
  Chieh-Hsin Lai, Yang Song, Dongjun Kim, Yuki Mitsufuji,
  Stefano Ermon,\\
  \emph{The Principles of Diffusion Models}, arXiv:2510.21890.\\[0.3em]
  \href{https://arxiv.org/pdf/2510.21890}
       {\texttt{https://arxiv.org/pdf/2510.21890}}
\end{block}

\begin{itemize}
  \item 本早見表は同資料の「Notations」と第1章、および後続章で使う
        Score・SDE・ODEの記法を意識して構成
  \item 数式を引用する際は、元資料の式番号と定義を必ず再確認する
  \item コード例は必要な部分をコピーして各ゼミ資料へ利用できる
\end{itemize}
\end{frame}

\end{document}
